% 第三篇，保留原篇全部论证和公式。
\chapter{模 \texorpdfstring{$\ell^n$}{l n} 与模 \texorpdfstring{$p$}{p} 的 \texorpdfstring{$L$}{L} 函数}\label{III}
\origpage{116}{110}
本篇证明《迹公式报告》（本书，以下简写为“报告”）中定理的一种推广：系数环 $A$ 可以是与 $p$ 互素的挠环，也可以是特征 $p$ 的域。证明将实质性地使用 \cite[XVII 5.5]{sga4}。

\section{对称张量}\label{III:1}
\subsection{}\label{III:1-1}
设 $A$ 为交换环。关于函子 $\Gamma^n:(A\text{-模})\to(A\text{-模})$ 的性质，参见 \cite[XVII 5.5.1--2]{sga4}。这里只回顾：若 $M$ 平坦（实际只需考虑 $M$ 有限生成投射的情形），则
\[\Gamma^nM=(M^{\otimes n})^{S_n}\qquad\text{（$n$ 次对称张量）。}\]
这里称为\term{对称张量}{symmetric tensor}{tenseur symétrique}。

\subsection{}\label{III:1-2}
设 $X$ 为 $S$-概形，$X^n/S$ 为 $X$ 的 $n$ 次纤维积，$\pi$ 为从 $X^n/S$ 到 $(X^n/S)/S_n=\operatorname{Sym}_S^nX$ 的投影，并假定此商存在。若 $\sG$ 为 $X$ 上的 $A$-模层，则 $n$ 份 $\sG$ 的\term{外张量积}{external tensor product}{produit tensoriel externe} $\sG^{\boxtimes n}$ 是 $X^n/S$ 上的 $S_n$-等变层。我们将使用 \cite[XVII 5.5.7--9]{sga4} 的\emph{外} $\Gamma^n$ 函子。这里只回顾：若 $\sG$ 为平坦 $A$-模层（实际只需可构造平坦的情形），则
\[\Gamma_{\mathrm{ext}}^n(\sG)=(\pi_*\sG^{\boxtimes n})^{S_n}.\]

\begin{lemma_}\label{III:1-3}
设 $S=\spec k$，其中 $k$ 代数闭，$X$ 为有限份 $S$ 的不交并，$\sG$ 为 $X$ 上的平坦 $A$-模层。则
\[\Gamma(\operatorname{Sym}_S^nX,\Gamma_{\mathrm{ext}}^n\sG)=\Gamma^n\Gamma(X,\sG).\]
\end{lemma_}
事实上，
\begin{align*}
\Gamma(\operatorname{Sym}_S^nX,\Gamma_{\mathrm{ext}}^n\sG)
&=\Gamma(\operatorname{Sym}_S^nX,\pi_*\sG^{\boxtimes n})^{S_n}\\
&=\Gamma(X^n/S,\sG^{\boxtimes n})^{S_n}
=(\Gamma(X,\sG)^{\otimes n})^{S_n}=\Gamma^n\Gamma(X,\sG).
\end{align*}

\subsection{}\label{III:1-4}
以下需要非加性函子 $\Gamma^n$ 与 $\Gamma_{\mathrm{ext}}^n$ 的导出函子。这类导出函子的理论由 Dold 和 Puppe 建立 \cite{do61}。比以下回顾更完整的介绍，见 \cite[XVII 5.5.3]{sga4}。
设 $\cA$ 为阿贝尔范畴；甚至只需它是加性范畴，且每个幂等自身态射都是某个直和因子上的投影。以 $N$ 表示\term{正规化函子}{normalization functor}{foncteur de normalisation}：
\[N:(\cA\text{ 中的余单纯对象})\longrightarrow(\cA\text{ 中满足 $K^i=0$（$i<0$）的微分复形}).\]
这是范畴等价。它及其逆都将同伦态射变成同伦态射。若 $T:\cA\to\cB$ 为函子，$K$ 为 $\cA$ 中非负次数的复形，即 $i<0$ 时 $K^i=0$，则定义 $TK=NTN^{-1}K$。

\subsection{例}\label{III:1-5}
若 $K$ 仅在 $0$ 次为 $M$，其余次数为零，则 $TK$ 仅在 $0$ 次为 $TM$。

\subsection{}\label{III:1-6}
若 $K$ 是由有限生成投射 $A$-模组成的非负次数有界复形，则 $A$-模复形 $\Gamma^nK$ 也有这些性质。
以 $\eD_{\mathrm{parf}}^{\geqslant0}(A)$ 表示 $\eD_{\mathrm{parf}}(A)$（报告 \ref{II:4-3}）中 Tor 维数不超过 $0$ 的复形所成的子范畴。由上可知，$\Gamma^n$ 诱导函子
\[\eL\Gamma^n:\eD_{\mathrm{parf}}^{\geqslant0}(A)\longrightarrow\eD_{\mathrm{parf}}^{\geqslant0}(A).\]
$\eL\Gamma_{\mathrm{ext}}^n$ 的类似定义较为复杂，参见 \cite[XVII 5.5.14]{sga4}。

\subsection{}\label{III:1-7}
设 $K$ 为有限生成投射 $A$-模的有界复形，$u$ 为 $K$ 的自身态射。以 $\det(1-ut,K)$ 表示常数项为 $1$ 的形式幂级数
\[\det(1-ut,K)=\prod_i\det(1-ut,K^i)^{(-1)^i}.\]
以下性质留给读者验证。
\subsubsection{}\label{III:1-7-1}
设 $F$ 为 $u$ 保持的有限滤过，且各 $\gr_F^p(K^q)$ 有限生成投射。则
\[\det(1-ut,K)=\prod_p\det(1-ut,\gr_F^pK).\]
\subsubsection{}\label{III:1-7-2}
将 $K$ 向左平移 $q$ 格，有
\[\det(1-ut,K[q])=\det(1-ut,K)^{(-1)^q}.\]

\begin{proposition_}\label{III:1-8}
设 $u$ 为有限生成投射 $A$-模的非负次数有界复形 $K$ 的自身态射。则
\begin{equation*}\tag{1.8.1}\label{III:eq:1-8-1}
\det(1-ut,K)^{-1}=\sum_{n\geqslant0}\tr(u,\Gamma^nK)t^n.
\end{equation*}
\end{proposition_}
将 \eqref{III:eq:1-8-1} 右边记为 $D(u,K)$。
\begin{lemma_}\label{III:1-9}
设 $u$ 为有限生成投射 $A$-模的非负次数有界复形短正合列
\[0\to K'\to K\to K''\to0\]
的自身态射。则 $D(u,K)=D(u,K')D(u,K'')$。
\end{lemma_}
若 $0\to M'\to M\to M''\to0$ 是可分裂的 $A$-模短正合列，则 $\Gamma^nM$ 有一个自然滤过，其相继商为 $\Gamma^iM'\otimes\Gamma^jM''$（$i+j=n$）。因此 $\Gamma^nK$ 有自然滤过，其相继商为复形
\[N(N^{-1}\Gamma^iK'\otimes N^{-1}\Gamma^jK'')\qquad(i+j=n).\]
它们典范地同伦等价于双复形 $\Gamma^iK'\otimes\Gamma^jK''$ 的单复形 $s(\Gamma^iK'\otimes\Gamma^jK'')$。故
\[\tr(u,\Gamma^nK)=\sum_{i+j=n}\tr(u,\Gamma^iK')\tr(u,\Gamma^jK''),\]
引理得证。由 \ref{III:1-9} 归纳得到：
\begin{lemma_}\label{III:1-10}
设 $F$ 为 $K$ 上 $u$ 保持的有限滤过，且各 $\gr_F^p(K^q)$ 有限生成投射。则
\[D(u,K)=\prod_pD(u,\gr_F^pK).\]
\end{lemma_}

\begin{lemma_}\label{III:1-11}
设 $u$ 为有限生成投射 $A$-模 $M$ 的自身态射，$M[-q]$ 为仅在 $q$ 次为 $M$ 的复形。则对 $q\geqslant0$，有
\[D(u,M[-q])=\left(\sum_{n\geqslant0}\tr(u,\Gamma^nM)t^n\right)^{(-1)^q}.\]
\end{lemma_}
由 \ref{III:1-5}，$q=0$ 时成立。对 $q$ 归纳，只需证明
\[D(u,M[-q])D(u,M[-(q+1)])=1.\]
将 \ref{III:1-9} 用于仅在 $q$、$q+1$ 次为 $M$、微分为恒等映射的复形 $K$，以及它的朴素滤过
\[(\cdots\to0\to M\to\cdots)\longrightarrow(\cdots\to M\xrightarrow{\mathrm{id}}M\to\cdots)\longrightarrow(\cdots\to M\to0\to\cdots),\]
即可得此式。事实上，$K$ 零伦，因而在导出范畴中 $\Gamma^nK=\Gamma^n0$，即 $n>0$ 时为 $0$、$n=0$ 时为 $A$，故 $D(u,K)=1$。

\subsection{}\label{III:1-12}
现在证明 \ref{III:1-8}。将 \ref{III:1-10} 用于由子复形
\[0\to\cdots\to0\to K^q\to K^{q+1}\to\cdots\]
定义的朴素滤过，再用 \ref{III:1-11}，得
\[D(u,K)=\prod_qD(u,K^q[-q])=\prod_q\left(\sum_{n\geqslant0}\tr(u,\Gamma^nK^q)t^n\right)^{(-1)^q}.\]
因此只需对有限生成投射模 $M$ 的自身态射 $u$ 证明
\begin{equation*}\tag{1.12.1}\label{III:eq:1-12-1}
\det(1-ut,M)^{-1}=\sum_{n\geqslant0}\tr(u,\Gamma^nM)t^n.
\end{equation*}
若 $u=0$，则两边均为 $1$。给 $M$ 加上一个 $u$ 在其上为零的模 $N$，并用 \ref{III:1-9}，可设 $M$ 自由。选取一组基后，所需证明的就是关于 $u$ 的矩阵系数的代数恒等式。由代数恒等式的延拓原理，可设 $A$ 为代数闭域。此时 $M$ 有 $u$ 保持的滤过，且相继商秩为 $1$。由 \ref{III:1-10}，只需考虑 $M$ 秩为 $1$，此时 \eqref{III:eq:1-12-1} 化为
\[\frac1{1-at}=\sum_{n\geqslant0}a^nt^n.\]

\begin{corollary_}\label{III:1-13}
设 $u$、$u'$ 分别为有限生成投射 $A$-模的有界复形 $K$、$K'$ 的自身态射。若在导出范畴中，带 $u$ 的 $K$ 与带 $u'$ 的 $K'$ 同构，则
\[\det(1-ut,K)=\det(1-u't,K').\]
\end{corollary_}
应用 \ref{III:1-7-2}，可设 $K$、$K'$ 都集中在非负次数。此时 \ref{III:1-8} 右边显然具有所需的不变性。

\subsection{}\label{III:1-14}
由此推论，可对 $K\in\ob\eD_{\mathrm{parf}}(A)$ 的自身态射 $u$ 定义 $\det(1-ut,K)$。这一构造是专门为当前目的设计的；实际上，对 $L\in\ob\eD_{\mathrm{parf}}(\Lambda)$ 的自同构 $v$，可定义 $\det v\in\Lambda^\times$。取 $v=1-ut$、$\Lambda=A\llbracket t\rrbracket$、$L=K\otimes_A\Lambda$，就得到 $\det(1-ut,K)$。

\section{定理}\label{III:2}
\subsection{}\label{III:2-1}
沿用报告第 \ref{II:1} 节的记号。设 $A$ 为交换诺特挠环，$X_0$ 为 $\dF_q$ 上的分离有限型概形，$\sK_0\in\ob\eD_{\mathrm{ctf}}(X_0;A)$。定义
\[L(X_0,\sK_0)=\prod_{x\in|X_0|}\det(1-F_x^*t^{\deg(x)},\sK)^{-1}\]
（参见 \ref{III:1-7}、\ref{III:1-13} 及报告 \ref{II:1-5}、\ref{II:1-6}）。

\begin{theorem_}\label{III:2-2}
在下述任一情形下：
\begin{enumerate}[\indent a)]
\item $A$ 为与 $p$ 互素的挠环；
\item $A$ 约化且特征为 $p$；
\end{enumerate}
都有
\begin{equation*}\tag{2.2.1}\label{III:eq:2-2-1}
L(X_0,\sK_0)=\det(1-F^*t^f,\eR\Gamma_c(X,\sK))^{-1}.
\end{equation*}
\end{theorem_}

\subsection{注}\label{III:2-3}
在 a) 情形，报告第 \ref{II:3} 节的方法可从迹公式（报告 \ref{II:4-10}）推出 \eqref{III:eq:2-2-1} 两边的对数导数相等。但由于 $A$ 为挠环，这还不足以证明 \ref{III:2-2}。

\subsection{注}\label{III:2-4}
若 $A$ 为域，则
\[L(X_0,\sK_0)=\prod_iL(X_0,\underline\h^i(\sK_0))^{(-1)^i}.\]
谱序列 $E_2^{pq}=\h_c^p(X,\underline\h^q(\sK))\Longrightarrow\h_c^{p+q}(X,\sK)$ 给出
\begin{align*}
\det(1-F^*t^f,\eR\Gamma_c(X,\sK))
&=\prod_n\det(1-F^*t^f,\h_c^n(X,\sK))^{(-1)^n}\\
&=\prod_{p,q}\det(1-F^*t^f,\h_c^p(X,\underline\h^q(\sK)))^{(-1)^{p+q}}.
\end{align*}
这将 $A$ 为域时的 \ref{III:2-2} 归约到以下特殊情形。
\subsubsection{}\label{III:2-4-1}
若 $\sG_0$ 为可构造 $A$-向量空间层，则
\[L(X_0,\sG_0)=\prod_n\det(1-F^*t^f,\h_c^n(X,\sG))^{(-1)^{n+1}}.\]

\subsection{注}\label{III:2-5}
b) 情形容易归约到 $A$ 是特征 $p$ 的域、甚至是有限域的情形，参见 \ref{III:4-2}。
\subsection{注}\label{III:2-6}
b) 中 $A=\dZ/p$、$\sK$ 仅在 $0$ 次为常值层 $\dZ/p$ 的特殊情形，即
\[Z(X_0,t)=\prod_{x\in|X_0|}(1-t^{\deg(x)})^{-1}
\equiv\prod_n\det(1-F^*t^f,\h_c^n(X,\dZ/p))^{(-1)^{n+1}}\pmod p,\]
已由 N. Katz 证明 \cite[XXII 3.1]{sga7}。他从 Dwork 关于 $Z(X_0,t)$ 的 $p$-进理论出发，所用方法与这里完全不同。
\subsection{注}\label{III:2-7}
与我先前贸然向朋友们作出的断言相反，b) 中不能删去“$A$ 约化”的假设。反例见 \ref{III:4-5}。

\subsection{\texorpdfstring{\ref{III:2-2}}{2.2} 的证明纲要}\label{III:2-8}
由次数平移 \ref{III:1-7-2}，可设 $\sK_0$ 是非负次数、逐项平坦可构造的有界 $A$-模层复形。此时 $\eR\Gamma_c(X,\sK)\in\ob\eD_{\mathrm{parf}}^{\geqslant0}(A)$（报告 \ref{II:4-9} 的证明 b)）。我们将用 \ref{III:1-8} 把 \eqref{III:eq:2-2-1} 两边展开为 $T=t^f$ 的幂级数。$T^n$ 系数的相等将由 \cite[XVII 5.5.21]{sga4} 及 $\operatorname{Sym}^nX$ 上的迹公式推出：a) 情形用报告 \ref{II:4-10}，b) 情形用下文将证明的公式。
首先应归约到 $X_0$ 的对称幂存在的情形，例如 $X_0$ 仿射。为此使用 \eqref{III:eq:2-2-1} 两边关于 $X_0$ 的乘法性：若 $U_0$ 为 $X_0$ 的开子集，补闭集为 $Y_0$，则
\[L(X_0,\sK_0)=L(U_0,\sK_0|U_0)L(Y_0,\sK_0|Y_0).\]
右边行列式的相应公式由 \ref{III:1-7-1} 按报告第 \ref{II:6} 节（A）的方法推出。

\subsection{\texorpdfstring{$X_0$}{X0} 有限的情形}\label{III:2-9}
\eqref{III:eq:2-2-1} 两边关于 $X_0$ 的乘法性将问题归约到 $X_0=\spec(\dF_{q^k})$。回到定义可知，此时 \eqref{III:eq:2-2-1} 等价于报告 \ref{II:3-4}。

\subsection{左边}\label{III:2-10}
由 \ref{III:1-8}，有
\[L(X_0,\sK_0)=\prod_{x\in|X_0|}\sum_{m\geqslant0}\tr(F_x^*,\Gamma^m\sK)t^{m\deg(x)}.\]
因此 $T^n$ 的系数，是对所有满足 $\sum_xn_x[k(x):\dF_q]=n$ 的形式线性组合 $\sum_{x\in|X_0|}n_xx$（$n_x\geqslant0$，且仅有限多个非零），将下列乘积求和：
\begin{equation*}\tag{2.10.1}\label{III:eq:2-10-1}
\prod_x\tr(F_x^*,\Gamma^{n_x}\sK).
\end{equation*}
对 $x\in|X_0|$，以 $(x)$ 表示 $X$ 中位于 $x$ 上方的 $[k(x):\dF_q]$ 个点组成的零维循环，它们构成 $F$ 的一个轨道。映射 $\sum_xn_xx\mapsto\sum_xn_x(x)$，将上述满足次数条件的形式组合，与 $X$ 上次数为 $n$ 且被 $F$ 固定的有效零维循环等同，也就是与 $\operatorname{Sym}^nX$ 中 $F$ 的不动点等同。因此 $L(X_0,\sK_0)$ 中 $T^n$ 的系数，是以 $(\operatorname{Sym}^nX)^F$ 为指标的乘积 \eqref{III:eq:2-10-1} 之和。更确切地说：
\begin{lemma_}\label{III:2-11}
$L(X_0,\sK_0)=\sum_{n\geqslant0}\sum_{x\in(\operatorname{Sym}^nX)^F}\tr(F_x^*,\eL\Gamma_{\mathrm{ext}}^n\sK)T^n$。
\end{lemma_}
令 $Y_0$ 遍及 $X_0$ 中越来越大的有限子概形，则
\[L(X_0,\sK_0)=\lim_{Y_0}L(Y_0,\sK_0|Y_0).\]
也就是说，对每个 $m$，存在有限子概形 $Z_0\subset X_0$，使只要 $Y_0\supset Z_0$，$L(X_0,\sK_0)$ 与 $L(Y_0,\sK_0|Y_0)$ 就模 $T^m$ 同余。只需让 $Z_0$ 包含次数不超过 $m$ 的所有闭点。右边也有同样性质，因此只需在 $X_0$ 有限时证明。此时由 \ref{III:2-9}、\ref{III:1-8}、\ref{III:1-3}，依次有
\begin{align*}
L(X_0,\sK_0)&=\det(1-F^*T,\Gamma(X,\sK))^{-1}\\
&=\sum_{n\geqslant0}\tr(F^*,\Gamma^n\Gamma(X,\sK))T^n\\
&=\sum_{n\geqslant0}\tr(F^*,\Gamma(\operatorname{Sym}^nX,\Gamma_{\mathrm{ext}}^n\sK))T^n\\
&=\sum_{n\geqslant0}\sum_{x\in(\operatorname{Sym}^nX)^F}\tr(F_x^*,\Gamma_{\mathrm{ext}}^n\sK)T^n.
\end{align*}
最后一步用了有限集合这一平凡情形的迹公式。

\subsection{右边}\label{III:2-12}
由 \ref{III:1-8} 及 \cite[XVII 5.5.21]{sga4}，有
\begin{align*}
\det(1-F^*T,\eR\Gamma_c(X,\sK))^{-1}
&=\sum_{n\geqslant0}\tr(F^*,\eL\Gamma^n\eR\Gamma_c(X,\sK))T^n\\
&=\sum_{n\geqslant0}\tr(F^*,\eR\Gamma_c(\operatorname{Sym}^nX,\eL\Gamma_{\mathrm{ext}}^n\sK))T^n.
\end{align*}
\subsection{}\label{III:2-13}
故定理等价于以下迹公式：
\[\tr(F^*,\eR\Gamma_c(\operatorname{Sym}^nX,\eL\Gamma_{\mathrm{ext}}^n\sK))
=\sum_{x\in(\operatorname{Sym}^nX)^F}\tr(F_x^*,\eL\Gamma_{\mathrm{ext}}^n\sK).\]
a) 情形由报告 \ref{II:4-10} 得到；b) 情形的相应公式将在第 \ref{III:4} 节证明。

\section{Artin--Schreier 理论}\label{III:3}
\subsection{}\label{III:3-1}
以下总将概形 $S$ 上的凝聚层视为平展景 $S_{\mathrm{\acute et}}$ 上的层。设 $S$ 的特征为 $p>0$。若 $G$ 为有限秩局部常值 $\dZ/p$-向量空间层，则 $\sG=G\otimes_{\dZ/p}\sO$ 为 $S$ 上的局部自由凝聚层。以 $\Phi$ 同时表示 $\sO$ 的自身态射 $f\mapsto f^p$，及它与 $G$ 上恒等映射的张量积 $\Phi:\sG\to\sG$。Artin--Schreier 理论 \cite[IX 3.5]{sga4} 断言下列层序列正合：
\[\xymatrix{0\ar[r]&\dZ/p\ar[r]&\sO\ar[r]^-{\Phi-1}&\sO\ar[r]&0}.\]
与 $G$ 作张量积，得到正合列
\[\xymatrix{0\ar[r]&G\ar[r]&\sG\ar[r]^-{\Phi-1}&\sG\ar[r]&0}.\]

\subsection{}\label{III:3-2}
设 $S$ 为诺特概形 $\bar S$ 的开子集，$j:S\hookrightarrow\bar S$ 为包含态射，$I$ 为定义补闭集 $\bar S\setminus S$ 的理想层。将 $\Phi$ 经函子性延拓到 $j_*\sG$ 后，仍记为 $\Phi$。
\begin{lemma_}\label{III:3-3}
存在 $\sG$ 的凝聚延拓 $\bar\sG\subset j_*\sG$，使
\begin{equation*}\tag{3.3.1}\label{III:eq:3-3-1}
\Phi(\bar\sG)\subset I\bar\sG.
\end{equation*}
对任一这样的延拓，有正合列
\begin{equation*}\tag{3.3.2}\label{III:eq:3-3-2}
\xymatrix{0\ar[r]&j_!G\ar[r]&\bar\sG\ar[r]^-{\Phi-1}&\bar\sG\ar[r]&0}.
\end{equation*}
\end{lemma_}
\begin{enumerate}[a)]
\item \emph{存在性。}取 $\sG$ 在 $\bar S$ 上的凝聚延拓 $\sG'\subset j_*\sG$。对所有 $n$，有 $\Phi(I^n\sG')\subset I^{pn}\Phi(\sG')$。取 $n$ 充分大，使 $I^{n(p-1)-1}\Phi(\sG')\subset\sG'$，并置 $\bar\sG=I^n\sG'$。于是
\[\Phi(\bar\sG)\subset I^{pn}\Phi(\sG')=I\cdot I^n\cdot I^{n(p-1)-1}\Phi(\sG')\subset I\cdot I^n\sG'=I\bar\sG.\]
\item \emph{$\Phi-1$ 的核。}在 $S$ 上由 \ref{III:3-1} 得到结论。只需验证：对 $\bar S$ 上的平展 $U$，$\ker(\Phi-1)$ 在 $U$ 上的截面 $x$，在 $\bar S\setminus S$ 的逆像的某个邻域上为零。因为 $x=\Phi(x)$，由 \eqref{III:eq:3-3-1} 得 $x\in I\bar\sG$；又有
\[x\in I^n\bar\sG\quad\Longrightarrow\quad x=\Phi(x)\in\Phi(I^n\bar\sG)\subset I^{np}\bar\sG.\]
所以 $x$ 属于所有 $I^n\bar\sG$，从而在 $\bar S\setminus S$ 附近为零。
\item \emph{$\Phi-1$ 的满射性。}可设 $\bar S$ 仿射，并将 $\bar\sG$ 写成 $\sO^n$ 的商。此时 $\Phi$ 可提升为 $p$-线性映射 $\widetilde\Phi:\sO^n\to\sO^n$，只需证明 $\widetilde\Phi-1$ 满射：
\[\xymatrix{\sO^n\ar[r]^-{\widetilde\Phi-1}\ar[d]&\sO^n\ar[d]\\\bar\sG\ar[r]^-{\Phi-1}&\bar\sG}.\]
层 $\sO^n$ 是群概形 $\dG_a^n$ 的局部截面层，$\widetilde\Phi-1$ 由 $\dG_a^n\to\dG_a^n$ 的群同态诱导。这个同态平展。要验证其满射性，只需在每个代数闭域值点 $\spec k\to\bar S$ 上作基变换后验证；在代数闭域上，同维连通群之间的平展同态必为满射。最后，$\bar S$-概形之间的平展满态射在 $\bar S_{\mathrm{\acute et}}$ 上的相应层之间诱导满态射。
\end{enumerate}

\subsection{两种 Frobenius}\label{III:3-4}
\term{绝对 Frobenius 态射}{absolute Frobenius morphism}{morphisme de Frobenius absolu} $F_{\mathrm{abs}}:S\to S$ 在 $S$ 的底拓扑空间上为恒等映射，在结构层上为 $f\mapsto f^p$。对平展 $U/S$，有典范同构 $F_{\mathrm{abs}}^*U=U$，即下图为笛卡儿图：
\[\xymatrix{U\ar[r]^-{F_{\mathrm{abs}}}\ar[d]&U\ar[d]\\S\ar[r]^-{F_{\mathrm{abs}}}&S}.\]
因此，$F_{\mathrm{abs}}$ 定义的环化景自身态射 $(S_{\mathrm{\acute et}},\sO)\to(S_{\mathrm{\acute et}},\sO)$ 可描述为：
\begin{center}
$F_{\mathrm{abs}}^*$ 在 $S_{\mathrm{\acute et}}$ 上为恒等函子；\par
同态 $F_{\mathrm{abs}}^*\sO=\sO\to\sO$ 为 $\Phi$。
\end{center}
在此描述下，对 $S$ 上任一平展层 $G$，Frobenius 对应 $F_{\mathrm{abs}}^*G\to G$ 就是恒等映射。
给定 $S$ 上拟凝聚层的 $p$-线性态射 $u:\sM\to\sN$，等价于给定线性态射 $u':F_{\mathrm{abs}}^*\sM\to\sN$。特别地，沿用以上记号，同态 $\Phi:\sG\to\sG$ 对应于 $\Phi':F_{\mathrm{abs}}^*\sG\to\sG$。

\begin{lemma_}\label{III:3-5}
沿用 \ref{III:3-4} 的记号，下图交换：
\[\xymatrix{F_{\mathrm{abs}}^*G\ar[r]\ar[d]&G\ar[d]\\F_{\mathrm{abs}}^*\sG\ar[r]^-{\Phi'}&\sG}.\]
\end{lemma_}
只需局部验证，所以可设 $G$ 常值，此时显然。$F_{\mathrm{abs}}$ 的各次幂也有同样的理论。

\subsection{}\label{III:3-6}
现在设 $\bar S$ 为 $\dF_q$（$q=p^f$）上的概形；按报告第 \ref{II:1} 节的约定，将 $S,\bar S,G,\sG,\bar\sG$ 改记为 $S_0,\bar S_0,G_0,\sG_0,\bar\sG_0$。将引理 \ref{III:3-5} 用于 $F_{\mathrm{abs}}^f=F:S_0\to S_0$，得到交换图
\[\xymatrix{F^*G_0\ar[r]\ar[d]&G_0\ar[d]\\F^*\sG_0\ar[r]^-{(\Phi^f)'}&\sG_0},\qquad
\text{进而有}\qquad
\xymatrix{F^*(j_!G_0)\ar[r]\ar[d]&j_!G_0\ar[d]\\F^*\bar\sG_0\ar[r]^-{(\Phi^f)'}&\bar\sG_0}.\]
再从 $\dF_q$ 向 $\dF$ 扩张标量，并取上同调，得到
\[\xymatrix{F^*(j_!G)\ar[r]^-{F^*}\ar[d]&j_!G\ar[d]\\F^*\bar\sG\ar[r]&\bar\sG}
\qquad\text{及}\qquad
\xymatrix{\h^\bullet(\bar S,j_!G)\ar[r]^-{F^*}\ar[d]&\h^\bullet(\bar S,j_!G)\ar[d]\\\h^\bullet(\bar S,\bar\sG)\ar[r]&\h^\bullet(\bar S,\bar\sG)}.\]
这里下方的作用由 $\h^\bullet(\bar S_0,\bar\sG_0)$ 上的 $\dF_q$-线性自身态射 $\Phi^f$ 经 $\dF_q\to\dF$ 的标量扩张得到。
Artin--Schreier 理论给出长正合列
\[\xymatrix{\cdots\ar[r]&\h^i(\bar S,j_!G)\ar[r]&\h^i(\bar S,\bar\sG)\ar[r]^-{\Phi-1}&\h^i(\bar S,\bar\sG)\ar[r]&\cdots}.\]
其中 $\Phi$ 是 $\h^\bullet(\bar S_0,\bar\sG_0)$ 的 $p$-线性自身态射 $\Phi$ 在
\[\h^\bullet(\bar S,\bar\sG)=\h^\bullet(\bar S_0,\bar\sG_0)\otimes_{\dF_q}\dF\]
上的 $p$-线性扩张。

\subsection{}\label{III:3-7}
若 $\bar S_0$ 在 $\dF_q$ 上固有，则各 $\h^i(\bar S,\bar\sG)$ 有限维，且 $\Phi-1$ 满射，证明同 \ref{III:3-3} c)。故
\[\h_c^i(S,G)=\ker\bigl(\Phi-1:\h^i(\bar S,\bar\sG)\to\h^i(\bar S,\bar\sG)\bigr),\]
而 $F^*$ 由 $\h^i(\bar S_0,\bar\sG_0)$ 的 $\dF_q$-线性自身态射 $\Phi^f$，在 $\h^i(\bar S,\bar\sG)=\h^i(\bar S_0,\bar\sG_0)\otimes_{\dF_q}\dF$ 上的 $\dF$-线性扩张诱导。
应用下面的引理 \ref{III:3-8}，得到等式
\begin{equation*}\tag{3.7.1}\label{III:eq:3-7-1}
\tr(F^*,\h_c^i(S,G))=\tr(\Phi^f,\h^i(\bar S_0,\bar\sG_0)).
\end{equation*}

\begin{lemma_}\label{III:3-8}
设 $\Phi$ 为 $\dF_q$ 上有限维向量空间 $V$ 的 $p$-线性自身态射，即 $\Phi(\lambda x)=\lambda^p\Phi(x)$。置 $F=\Phi^f$，这是线性映射；将 $\Phi$、$F$ 在 $V\otimes_{\dF_q}\dF$ 上分别作 $p$-线性、线性扩张，仍记为 $\Phi$、$F$。则 $F$ 保持 $\ker(\Phi-1)\subset V\otimes_{\dF_q}\dF$，且
\[\tr(F,\ker(\Phi-1))=\tr(F,V)=\tr(F,V\otimes_{\dF_q}\dF).\]
\end{lemma_}
\translatornote{原文未明确写出 $V$ 有限维；这里补入这一条件，才能使用通常的迹与下文的 Fitting 分解，本节应用已满足它。另外，原书第124页证明的第一个张量积将 $\ker(\Phi-1)\otimes\dF$ 的底域印为 $\dF_q$，应为 $\dF_p$，下文的迹计算已经使用了正确的 $\dF_p$。确实，若 $\Phi(x)=x$，则 $\Phi(\lambda x)=\lambda^px$，故不动向量构成 $\dF_p$-向量空间，通常并不构成 $\dF_q$-向量空间；在 $\Phi$ 可逆的部分，半线性下降给出其经 $\dF_p\to\dF$ 扩张后的同构。}
在 $V\otimes_{\dF_q}\dF$ 上，有 $F\Phi=\Phi F$，因为 $p$-线性映射 $F\Phi-\Phi F$ 在 $V$ 上为零，因而处处为零。将 $V$ 分解为 $V'\oplus V''$，其中 $F$ 在 $V'$ 上可逆、在 $V''$ 上幂零。$p$-线性自身态射理论给出
\[\ker(\Phi-1)\otimes_{\dF_p}\dF\iso V'\otimes_{\dF_q}\dF.\]
由于 $\tr(F,V'')=0$，有
\begin{align*}
\tr(F,\ker(\Phi-1))&=\tr(F,\ker(\Phi-1)\otimes_{\dF_p}\dF)\\
&=\tr(F,V'\otimes_{\dF_q}\dF)\\
&=\tr(F,V\otimes_{\dF_q}\dF)=\tr(F,V).
\end{align*}

\section{模 \texorpdfstring{$p$}{p} 的迹公式}\label{III:4}
\begin{theorem_}\label{III:4-1}
设 $X_0$ 为 $\dF_q$ 上的分离有限型概形，$A$ 为特征 $p$ 的约化交换诺特环，$\sK_0\in\ob\eD_{\mathrm{ctf}}(X_0,A)$。则
\begin{equation*}\tag{4.1.1}\label{III:eq:4-1-1}
\sum_{x\in X^F}\tr(F_x^*,\sK_x)=\tr(F^*,\eR\Gamma_c(X,\sK)).
\end{equation*}
\end{theorem_}

\subsection{归约到 \texorpdfstring{$A$}{A} 为有限域的情形}\label{III:4-2}
标准的取极限论证将问题归约到 $A$ 是有限生成 $\dZ/p$-代数的情形。对 $A$ 的每个极大理想 $\fm$，\eqref{III:eq:4-1-1} 两边在 $A/\fm$ 中的像由同一公式给出，只需以 $\sK_0\lotimes_A A/\fm\in\ob\eD_{\mathrm{ctf}}(X_0,A/\fm)$ 替换 $\sK_0$；右边的相容性见报告 \ref{II:4-12}。各 $A/\fm$ 都是有限域，且 $A$ 的所有极大理想之交为 $0$，归约得证。

\subsection{归约到 \texorpdfstring{$A=\dZ/p$}{A=Z/p} 的情形}\label{III:4-3}
设 $A$ 为有限域，以 $T'_A(\sK_0)$、$T''_A(\sK_0)$ 表示 \eqref{III:eq:4-1-1} 的两边。通过限制标量，将 $\sK_0$ 看作 $\dZ/p$-模层的复形，则
\[T'_{\dZ/p}(\sK_0)=\tr_{A/(\dZ/p)}T'_A(\sK_0),\]
$T''$ 也同样。因此，若 \eqref{III:eq:4-1-1} 对 $A=\dZ/p$ 成立，则
\begin{equation*}\tag{4.3.1}\label{III:eq:4-3-1}
\tr_{A/(\dZ/p)}T'_A(\sK_0)=\tr_{A/(\dZ/p)}T''_A(\sK_0).
\end{equation*}
我断言，对任意 $\lambda\in A$，甚至有
\begin{equation*}\tag{4.3.2}\label{III:eq:4-3-2}
\tr_{A/(\dZ/p)}(\lambda T'_A(\sK_0))=\tr_{A/(\dZ/p)}(\lambda T''_A(\sK_0)),
\end{equation*}
从而 $T'_A(\sK_0)=T''_A(\sK_0)$。$\lambda=0$ 时显然；对可逆 $\lambda$，以 $A(\lambda)$ 表示 $\spec(\dF_q)$ 上满足 $F^*=\lambda$ 的秩为 $1$ 的自由 $A$-模层在 $X_0$ 上的逆像。\eqref{III:eq:4-3-2} 就是将 \eqref{III:eq:4-3-1} 用于 $A(\lambda)\otimes_A\sK_0$。

\subsection{}\label{III:4-4}
报告第 \ref{II:6} 节的方法现在将 \ref{III:4-1} 归约到下述陈述。
\begin{theorem*}[4.1 bis]\label{III:4-1bis}
设 $\bar S_0$ 为光滑射影曲线，$j:S_0\hookrightarrow\bar S_0$ 为稠密开子集，$G_0$ 为 $S_0$ 上的有限秩局部常值 $\dZ/p$-向量空间层。则
\begin{equation*}\tag{4.4.1}\label{III:eq:4-4-1}
\sum_{x\in S^F}\tr(F_x^*,G_x)=\sum_i(-1)^i\tr(F^*,\h_c^i(S,G)).
\end{equation*}
\end{theorem*}
\begin{proof}
将 \ref{III:3-3} 用于 $(S_0,\bar S_0,G_0)$。该引理给出的 $\sG_0$ 的凝聚延拓 $\bar\sG_0$ 局部自由：因为 $\bar\sG_0\subset j_*\sG_0$，它无挠，而 $\bar S_0$ 正则且一维。$\bar\sG_0$ 的 $q$-线性自身态射 $\Phi^f$ 可解释为态射
\[u:F^*\bar\sG_0\longrightarrow\bar\sG_0.\]
它通过 $I\bar\sG_0$ 分解，其中 $I$ 定义补闭集 $\bar S_0\setminus S_0$。以 $v:j_!G_0\to\bar\sG_0$ 表示自然映射 $G_0\to\sG_0=G_0\otimes_{\dZ/p}\sO$ 的唯一延拓。下图交换：
\begin{equation*}\tag{4.4.2}\label{III:eq:4-4-2}
\xymatrix{F^*j_!G_0\ar[r]^-{F^*}\ar[d]^-{F^*v}&j_!G_0\ar[d]^-v\\F^*\bar\sG_0\ar[r]^-u&\bar\sG_0}.
\end{equation*}
最后，以 $uF^*$ 表示复合
\[\xymatrix{\h^\bullet(\bar S,\bar\sG)\ar[r]^-{F^*}&\h^\bullet(\bar S,F^*\bar\sG)\ar[r]^-u&\h^\bullet(\bar S,\bar\sG)}.\]
由 \eqref{III:eq:3-7-1}，有
\begin{equation*}\tag{4.4.3}\label{III:eq:4-4-3}
\tr(F^*,\h_c^i(S,G))=\tr(uF^*,\h^i(\bar S,\bar\sG)).
\end{equation*}
对 $\bar S$ 的每个闭点 $i_x:x\hookrightarrow\bar S$，置 $\bar\sG_x=i_x^*\bar\sG$。将凝聚上同调中的 Lefschetz 迹公式（“Woods Hole 不动点公式”）用于 $F$ 与 $u$，得
\begin{equation*}\tag{4.4.4}\label{III:eq:4-4-4}
\sum_i(-1)^i\tr(uF^*,\h^i(\bar S,\bar\sG))=\sum_{x\in\bar S^F}\frac{\tr(u_x,\bar\sG_x)}{\det(1-dF_x)}.
\end{equation*}
右边各项满足：
\begin{enumerate}[\indent a)]
\item 因 $dF=0$，所有分母均为 $1$；
\item 对 $x\in\bar S^F\setminus S^F$，$\bar\sG_x$ 的自身态射 $u_x$ 为零，故 $\tr(u_x,\bar\sG_x)=0$；
\item 对 $x\in S^F$，\eqref{III:eq:4-4-2} 给出
\[\tr(u_x,\bar\sG_x)=\tr(F_x^*,G_x).\]
\end{enumerate}
于是 \eqref{III:eq:4-4-1} 由 \eqref{III:eq:4-4-3}、\eqref{III:eq:4-4-4} 推出。
\end{proof}

\subsection{一个反例}\label{III:4-5}
下面的例子说明，\ref{III:4-1} 中不能删去 $A$ 约化的假设。
设 $Y_0$ 为 $\dF_2$ 上 Frobenius 特征值为 $\pm\sqrt{-2}$ 的椭圆曲线。它有 $3$ 个有理点，且是\term{超奇异椭圆曲线}{supersingular elliptic curve}{courbe elliptique supersingulière}。因此对合 $\sigma:x\mapsto-x$ 只有一个不动点，即零点 $0$。$Y_0$ 关于 $\sigma$ 的商为射影直线，所以 $Y_0\setminus\{0\}$ 关于 $\sigma$ 的商为仿射直线 $\dA_0^1$，而 $Y_0\setminus\{0\}$ 是 $\dA_0^1$ 的无分歧二重覆盖：
\[\pi:Y_0\setminus\{0\}\longrightarrow\dA_0^1,\qquad\pi=\pi\sigma.\]
层 $\pi_*\dZ/2$ 在二元群 $\{1,\sigma\}$ 的 $\dZ/2$-群代数 $A$ 上局部自由、秩为 $1$。这个群代数同构于 $\dZ/2$ 上的\term{对偶数代数}{dual number algebra}{algèbre des nombres duaux}。
$A$-模层 $\pi_*\dZ/2$ 就是所需反例：
\begin{enumerate}[\indent a)]
\item 有 $\h_c^\bullet(\dA^1,\pi_*\dZ/2)=\h_c^\bullet(Y\setminus\{0\},\dZ/2)$。因 $Y_0$ 超奇异，$i>0$ 时 $\h^i(Y,\dZ/2)=0$，而长正合列
\[\cdots\to\h_c^i(Y\setminus\{0\},\dZ/2)\to\h^i(Y,\dZ/2)\to\h^i(\{0\},\dZ/2)\to\cdots\]
说明 $\h_c^\bullet(Y\setminus\{0\},\dZ/2)=0$。故 $\h_c^\bullet(\dA^1,\pi_*\dZ/2)=0$，从而
\[\tr_A(F^*,\eR\Gamma_c(\dA^1,\pi_*\dZ/2))=0.\]
\item 因 $Y_0\setminus\{0\}$ 有两个有理点，Frobenius 代换在 $\dA_0^1$ 的一个有理点上方为恒等映射，在另一个点上方为 $\sigma$。所以
\[\sum_{x\in(\dA^1)^F}\tr_A(F^*,\pi_*\dZ/2)=1+\sigma,\]
而 $1+\sigma\ne0$。
\end{enumerate}
覆盖 $\pi$ 的方程为 $y^2-y=x^3$。
